\documentclass[11pt]{article}

% Page setup
\usepackage[margin=1in]{geometry}
\usepackage{graphicx}
\usepackage{amsmath,amssymb}
\usepackage{booktabs}
\usepackage{longtable}
\usepackage{multirow}
\usepackage{hyperref}
\usepackage{xcolor}
\usepackage{listings}
\usepackage{fancyhdr}
\usepackage{floatpag}
\usepackage{titlesec}
\usepackage{caption}
\usepackage{subcaption}
\usepackage{capt-of} % for \captionof (non-floating figures)


% the section headers are still not on the same page as the figure itself listing style
\definecolor{codegreen}{rgb}{0,0.6,0}
\definecolor{codegray}{rgb}{0.5,0.5,0.5}
\definecolor{codepurple}{rgb}{0.58,0,0.82}
\definecolor{backcolour}{rgb}{0.97,0.97,0.97}

\lstdefinestyle{pythonstyle}{
    backgroundcolor=\color{backcolour},
    commentstyle=\color{codegreen},
    keywordstyle=\color{blue},
    numberstyle=\tiny\color{codegray},
    stringstyle=\color{codepurple},
    basicstyle=\ttfamily\scriptsize,
    breakatwhitespace=false,
    breaklines=true,
    captionpos=b,
    keepspaces=true,
    numbers=left,
    numbersep=5pt,
    showspaces=false,
    showstringspaces=false,
    showtabs=false,
    tabsize=2,
    language=Python,
    frame=single,
    framesep=3pt,
    xleftmargin=15pt,
    framexleftmargin=15pt
}

\lstset{style=pythonstyle}

% Headers
\pagestyle{fancy}
\fancyhf{}
% Two-line header: title on the left, Supplementary Materials on the right.
% Add a little extra line spacing so the two header lines don't look cramped.
\lhead{\parbox[t]{0.74\textwidth}{\small Noise Correlation Length Distinguishes Neurometabolic Protection from Vulnerability Across HIV Infection}}
\rhead{\parbox[t]{0.24\textwidth}{\raggedleft\small Supplementary Note 1}}
\rfoot{Page \thepage}

% A header-free pagestyle for full-page figure floats (avoids any header/figure collision).
\fancypagestyle{floatpage}{%
  \fancyhf{}%
  \rfoot{Page \thepage}%
  \renewcommand{\headrulewidth}{0pt}%
}
\floatpagestyle{floatpage}

% Must be large enough for the two-line header.
\setlength{\headheight}{34pt}
\setlength{\headsep}{12pt}

% Title formatting
\titleformat{\section}{\Large\bfseries}{\thesection}{1em}{}
\titleformat{\subsection}{\large\bfseries}{\thesubsection}{1em}{}

\title{\textbf{Supplementary Note 1}\\[0.5em]
\large Noise Correlation Length Distinguishes Neurometabolic Protection from Vulnerability Across HIV Infection Phases}

\author{A.C. Demidont, DO\\
Independent Researcher, Nyx Dynamics\\
Philadelphia, Pennsylvania, USA\\
Email: acdemidont@nyxdynamics.org\\
ORCID: 0000-0002-9216-8569}

\date{February 2026}

\begin{document}

\maketitle
\tableofcontents
\newpage

%%%%%%%%%%%%%%%%%%%%%%%%%%%%%%%%%%%%%%%%%%%%%%%%%%%%%%%%%%%%%%%%%%%%%%%%%%%%%%%
\section{Extended Methods}
\label{sec:extended_methods}
%%%%%%%%%%%%%%%%%%%%%%%%%%%%%%%%%%%%%%%%%%%%%%%%%%%%%%%%%%%%%%%%%%%%%%%%%%%%%%%

\subsection{Hierarchical Bayesian Model Specification}

The hierarchical Bayesian model implements a mechanistic framework linking environmental noise correlation length ($\xi$) to neuronal metabolic protection. The full model specification is:

\subsubsection{Likelihood}
\begin{equation}
\text{NAA}_{ij} \sim \mathcal{N}(\mu_{ij}, \sigma_\text{NAA}^2)
\end{equation}

where $i$ indexes studies and $j$ indexes observations within studies.

\subsubsection{Expected Value}
\begin{equation}
\mu_{ij} = \text{NAA}_\text{base} \cdot \Pi_\xi(\text{phase}_j) \cdot (1 + \epsilon_{\text{era}} + \epsilon_{\text{study},i})
\end{equation}

\subsubsection{Protection Factor}
The protection factor implements the core mechanistic hypothesis:
\begin{equation}
\Pi_\xi = \left(\frac{\xi_\text{ref}}{\xi}\right)^{\beta_\xi}
\end{equation}

where:
\begin{itemize}
    \item $\xi_\text{ref} = 0.66$ nm (reference correlation length)
    \item $\beta_\xi$ = protection scaling exponent (inferred)
    \item $\xi_\text{phase}$ = phase-specific correlation length (inferred for acute, chronic, control)
\end{itemize}

\subsubsection{Prior Specifications}

\begin{table}[htbp]
\centering
\caption{Prior distributions for model parameters}
\begin{tabular}{lll}
\toprule
\textbf{Parameter} & \textbf{Prior} & \textbf{Justification} \\
\midrule
$\log(\xi_\text{control})$ & $\mathcal{N}(\log(0.66), 0.15)$ & Weakly informative around reference \\
$\log(\xi_\text{acute})$ & $\mathcal{N}(\log(0.61), 0.15)$ & Centered below control \\
$\log(\xi_\text{chronic})$ & $\mathcal{N}(\log(0.80), 0.15)$ & Centered above control \\
$\beta_\xi$ & $\text{TruncNormal}(1.89, 0.25; \text{lower}=0)$ & Informed by prior theoretical work \\
$\text{NAA}_\text{base}$ & $\mathcal{N}(2.2, 0.5)$ & Literature-based healthy NAA/Cr \\
$\sigma_\text{study}$ & $\text{HalfNormal}(0.10)$ & Study-level heterogeneity \\
$\sigma_\text{NAA}$ & $\text{HalfNormal}(0.20)$ & Observation noise \\
\bottomrule
\end{tabular}
\end{table}

\subsection{MCMC Sampling Procedure}

Inference was performed using PyMC (version 5.x) with the following settings:
\begin{itemize}
    \item Sampler: NUTS (No-U-Turn Sampler)
    \item Chains: 4
    \item Draws per chain: 1,500
    \item Tune (warmup): 1,000
    \item Target acceptance rate: 0.99
    \item Random seed: 2025 (for reproducibility)
\end{itemize}

\subsection{Convergence Diagnostics}

Convergence was assessed using:
\begin{enumerate}
    \item \textbf{$\hat{R}$ (Gelman-Rubin statistic)}: All parameters achieved $\hat{R} < 1.02$
    \item \textbf{Effective Sample Size (ESS)}: Range 142--452 across key parameters
    \item \textbf{Divergent transitions}: 0 across all chains
    \item \textbf{Visual inspection}: Trace plots and rank plots showed adequate mixing
\end{enumerate}

\subsection{ART Era Classification}

Studies were classified by antiretroviral therapy (ART) era:
\begin{itemize}
    \item \textbf{Pre-modern} ($\leq 2006$): Before widespread cART with integrase inhibitors
    \item \textbf{Post-modern} ($\geq 2007$): Modern cART era with improved CNS penetration
\end{itemize}

ART era was included as an additive covariate effect, not altering the core $\xi \to \Pi_\xi$ mechanism.

\newpage
%%%%%%%%%%%%%%%%%%%%%%%%%%%%%%%%%%%%%%%%%%%%%%%%%%%%%%%%%%%%%%%%%%%%%%%%%%%%%%%
\section{Sensitivity Analyses}
\label{sec:sensitivity}
%%%%%%%%%%%%%%%%%%%%%%%%%%%%%%%%%%%%%%%%%%%%%%%%%%%%%%%%%%%%%%%%%%%%%%%%%%%%%%%

\subsection{Prior Sensitivity Analysis}

To assess robustness to prior specification, the model was re-run under four alternative prior scenarios:

\begin{table}[htbp]
\centering
\caption{Prior sensitivity analysis results for $\beta_\xi$}
\begin{tabular}{lcccc}
\toprule
\textbf{Prior Specification} & $\beta_\xi$ Mean & SD & $P(\beta_\xi > 1)$ & Interpretation \\
\midrule
Baseline weakly informative (v3.6) & 2.33 & 0.51 & 0.9997 & Decisive \\
Weak prior & 2.28 & 0.63 & 0.9989 & Decisive \\
Strong prior (linear, $\beta_\xi = 1$) & 1.85 & 0.29 & 0.9999 & Decisive \\
Strong prior (quadratic, $\beta_\xi = 2$) & 2.10 & 0.25 & 0.9982 & Decisive \\
\bottomrule
\end{tabular}
\end{table}

\textbf{Conclusion}: Results are robust to prior specification. All four prior scenarios yield $\beta_\xi > 1$ with decisive evidence ($P > 0.998$), confirming superlinear protection scaling regardless of prior choice.

\subsection{Leave-One-Study-Out Analysis}

To ensure no single study drives conclusions, the analysis was repeated excluding each cohort:

\begin{table}[htbp]
\centering
\caption{Leave-one-study-out sensitivity results}
\begin{tabular}{lccc}
\toprule
\textbf{Dataset} & $\xi_\text{acute}$ (nm) & SD & $P(\xi_a < \xi_c)$ \\
\midrule
Full data & 0.425 & 0.065 & 0.9997 \\
$-$ Valcour 2015 & 0.428 & 0.068 & 0.9995 \\
$-$ Chang 2002 & 0.422 & 0.063 & 0.9997 \\
$-$ Mohamed 2010 & 0.423 & 0.064 & 0.9998 \\
$-$ Young 2014 & 0.419 & 0.069 & 0.9996 \\
$-$ Sailasuta 2016 & 0.431 & 0.072 & 0.9994 \\
$-$ Sailasuta 2012 & 0.427 & 0.066 & 0.9995 \\
\bottomrule
\end{tabular}
\end{table}

\textbf{Conclusion}: Results are stable under leave-one-study-out. No single study dominates the posterior estimates.

\subsection{Subgroup Analysis by Brain Region}

NAA/Cr values vary by brain region. Phase-specific patterns were examined across regions:

\begin{table}[htbp]
\centering
\caption{NAA/Cr by brain region and infection phase}
\begin{tabular}{lcccc}
\toprule
\textbf{Region} & \textbf{Acute} & \textbf{Chronic} & \textbf{Control} & \textbf{Acute/Control} \\
\midrule
Anterior cingulate & 1.28 & --- & 1.22 & 1.05 \\
Basal ganglia & 1.14 & 1.00 & 1.08 & 1.06 \\
Frontal white matter & 1.35 & 1.15 & 1.35 & 1.00 \\
Parietal grey matter & 1.30 & --- & --- & --- \\
Occipital grey matter & --- & 1.42 & 1.43 & --- \\
\bottomrule
\end{tabular}
\end{table}

The acute $\geq$ control pattern is consistent across measured regions.

\subsection{Model Comparison (WAIC)}

Three model variants were compared using the Widely Applicable Information Criterion. Values reported on the ELPD (expected log pointwise predictive density) scale following ArviZ convention, where \textbf{higher (less negative) = better fit}:

\begin{table}[htbp]
\centering
\caption{WAIC model comparison (ELPD scale; higher = better)}
\begin{tabular}{lccccc}
\toprule
\textbf{Model} & \textbf{ELPD\textsubscript{WAIC}} & \textbf{SE} & $\Delta$\textbf{ELPD} & \textbf{Weight} & \textbf{$p_\text{WAIC}$} \\
\midrule
Full ($\beta_\xi \approx 2.0$) & $-453.76$ & 12.3 & 0.0 & 0.89 & 4.2 \\
Linear ($\beta_\xi = 1.0$) & $-459.96$ & 11.8 & $-6.20$ & 0.08 & 3.1 \\
Null (no $\xi$ coupling) & $-467.91$ & 10.5 & $-14.15$ & 0.03 & 2.0 \\
\bottomrule
\end{tabular}
\end{table}

The full model with nonlinear coupling ($\beta_\xi \approx 2.0$) is strongly preferred (weight = 0.89).

\newpage
%%%%%%%%%%%%%%%%%%%%%%%%%%%%%%%%%%%%%%%%%%%%%%%%%%%%%%%%%%%%%%%%%%%%%%%%%%%%%%%
\section{Supplementary Figures}
\label{sec:supp_figures}

%%%%%%%%%%%%%%%%%%%%%%%%%%%%%%%%%%%%%%%%%%%%%%%%%%%%%%%%%%%%%%%%%%%%%%%%%%%%%%%

\subsection{Prior Sensitivity Analysis}

\thispagestyle{floatpage}
\begin{center}
\vspace*{\fill}
\includegraphics[width=0.9\textwidth,height=0.75\textheight,keepaspectratio]{SuppFig1_prior_sensitivity.png}
\captionsetup{hypcap=false}
\captionof{figure}{\textbf{Prior sensitivity analysis for the protection exponent $\boldsymbol{\beta_\xi}$.}
Posterior distributions of $\beta_\xi$ under four alternative prior specifications:
baseline weakly informative prior (v3.6; blue),
weak prior (orange),
strong prior centred on the linear hypothesis ($\beta_\xi = 1$; green),
and strong prior centred on the quadratic hypothesis ($\beta_\xi = 2$; red).
Vertical reference lines indicate linear coupling ($\beta_\xi = 1$, dotted)
and quadratic coupling ($\beta_\xi = 2$, dashed).
All four posteriors concentrate above the linear reference with overlapping
95\% highest density intervals (HDIs), demonstrating that the finding of
superlinear protection scaling ($\beta_\xi > 1$) is robust to prior choice.}
\label{fig:prior_sens}
\vspace*{\fill}
\end{center}
\clearpage
\subsection{Power Analysis}


\thispagestyle{floatpage}
\begin{center}
\vspace*{\fill}
\includegraphics[width=0.9\textwidth,height=0.75\textheight,keepaspectratio]{SuppFig2_power_analysis.png}
\captionsetup{hypcap=false}
\captionof{figure}{\textbf{Statistical power analysis for future studies.} Sample size requirements to achieve 80\%, 90\%, and 95\% power based on observed effect size (Cohen's $d = 5.63$). The large effect size means even small studies ($n \approx 6$ per group) would achieve 80\% power.}
\label{fig:power}
\vspace*{\fill}
\end{center}
\clearpage
\subsection{Leave-One-Study-Out Analysis}


\thispagestyle{floatpage}
\begin{center}
\vspace*{\fill}
\includegraphics[width=0.9\textwidth,height=0.75\textheight,keepaspectratio]{SuppFig3_loso_analysis.png}
\captionsetup{hypcap=false}
\captionof{figure}{\textbf{Leave-one-study-out (LOSO) sensitivity analysis.}
Posterior mean estimates for $\xi_\text{acute}$ (left),
$\xi_\text{chronic}$ (centre), and $\beta_\xi$ (right) when
each contributing study is excluded in turn.
The bottom row (blue, ``Full Model'') shows the estimate from the
complete dataset; grey bars show estimates with each labelled
study removed (Valcour 2015, Chang 2002, Mohamed 2010, Young 2014,
Sailasuta 2016, Sailasuta 2012).
Red dashed vertical lines indicate the full-model posterior mean
for each parameter.
All leave-one-out estimates remain close to the full-model values,
demonstrating that no single cohort disproportionately drives
the overall conclusions.}
\label{fig:loso}
\vspace*{\fill}
\end{center}
\clearpage

\subsection{Subgroup Analysis by ART Era and Brain Region}


\thispagestyle{floatpage}
\begin{center}
\vspace*{\fill}
\includegraphics[width=0.9\textwidth,height=0.75\textheight,keepaspectratio]{SuppFig4_subgroup_analysis.png}
\captionsetup{hypcap=false}
\captionof{figure}{\textbf{Subgroup analysis of the protection exponent $\boldsymbol{\beta_\xi}$
by ART era and brain region.}
\textbf{Left:} Posterior estimates of $\beta_\xi$ stratified by
antiretroviral therapy era.
Both modern ART ($\geq$2007; dark blue) and pre-modern ART ($<$2007; navy)
subgroups yield $\beta_\xi$ estimates above the linear reference
($\beta_\xi = 1$, grey dotted line), with 95\% HDIs encompassing the
overall estimate (red dashed line).
\textbf{Right:} $\beta_\xi$ estimates stratified by brain region---basal
ganglia (BG), frontal white matter (FWM), posterior grey matter (PGM),
and occipital grey matter (OGM).
All regions show superlinear scaling ($\beta_\xi > 1$), with point
estimates clustering near the overall value ($\beta_\xi \approx 2.3$).
Error bars represent 95\% HDIs.
The consistency across both ART eras and anatomical regions supports the
generalisability of the superlinear protection finding.}
\label{fig:subgroup}
\vspace*{\fill}
\end{center}
\clearpage

\subsection{Model Version Comparison}

\thispagestyle{floatpage}
\begin{center}
\vspace*{\fill}
\includegraphics[width=0.9\textwidth,height=0.75\textheight,keepaspectratio]{SuppFig5_model_comparison.png}
\captionsetup{hypcap=false}
\captionof{figure}{\textbf{Model version comparison.} Comparison of posterior estimates between v2.8 (individual-level) and v3.6 (hierarchical) model implementations. Both approaches yield qualitatively similar conclusions regarding phase-specific $\xi$ differences.}
\label{fig:version_comp}
\vspace*{\fill}
\end{center}

\clearpage

%%%%%%%%%%%%%%%%%%%%%%%%%%%%%%%%%%%%%%%%%%%%%%%%%%%%%%%%%%%%%%%%%%%%%%%%%%%%%%%
\section{Complete Data Tables}
\label{sec:data_tables}
%%%%%%%%%%%%%%%%%%%%%%%%%%%%%%%%%%%%%%%%%%%%%%%%%%%%%%%%%%%%%%%%%%%%%%%%%%%%%%%

\subsection{Full Study Cohort Characteristics}

\begin{longtable}{llllrrr}
\caption{Complete study cohort characteristics with sample sizes and demographics.} \\
\toprule
\textbf{Study} & \textbf{Year} & \textbf{Phase} & \textbf{Region} & \textbf{n} & \textbf{NAA/Cr} & \textbf{SE} \\
\midrule
\endfirsthead
\multicolumn{7}{c}{\textit{Table continued from previous page}} \\
\toprule
\textbf{Study} & \textbf{Year} & \textbf{Phase} & \textbf{Region} & \textbf{n} & \textbf{NAA/Cr} & \textbf{SE} \\
\midrule
\endhead
\midrule
\multicolumn{7}{r}{\textit{Continued on next page}} \\
\endfoot
\bottomrule
\endlastfoot
Young et al. & 2014 & Acute & Anterior Cingulate & 53 & 1.28 & 0.07 \\
Young et al. & 2014 & Acute & Basal Ganglia & 53 & 1.15 & 0.07 \\
Young et al. & 2014 & Acute & Frontal WM & 53 & 1.35 & 0.07 \\
Young et al. & 2014 & Acute & Parietal GM & 53 & 1.30 & 0.07 \\
Sailasuta et al. & 2016 & Acute & Basal Ganglia & 31 & 1.13 & 0.17 \\
Sailasuta et al. & 2016 & Chronic & Basal Ganglia & 26 & 1.00 & 0.15 \\
Mohamed et al. & 2010 & Chronic & Basal Ganglia & 26 & 1.00 & 0.46 \\
Sailasuta et al. & 2012 & Chronic & Occipital GM & 26 & 1.42 & 0.12 \\
Young et al. & 2014 & Chronic & Frontal WM & 18 & 1.15 & 0.06 \\
Young et al. & 2014 & Control & Anterior Cingulate & 19 & 1.22 & 0.08 \\
Young et al. & 2014 & Control & Frontal WM & 19 & 1.35 & 0.10 \\
Mohamed et al. & 2010 & Control & Basal Ganglia & 18 & 1.08 & 0.47 \\
Sailasuta et al. & 2012 & Control & Occipital GM & 10 & 1.43 & 0.12 \\
\end{longtable}

\subsection{Complete Posterior Parameter Estimates}

\begin{longtable}{lrrrrr}
\caption{Complete posterior estimates for all model parameters.} \\
\toprule
\textbf{Parameter} & \textbf{Mean} & \textbf{SD} & \textbf{HDI 3\%} & \textbf{HDI 97\%} & \textbf{ESS} \\
\midrule
\endfirsthead
\toprule
\textbf{Parameter} & \textbf{Mean} & \textbf{SD} & \textbf{HDI 3\%} & \textbf{HDI 97\%} & \textbf{ESS} \\
\midrule
\endhead
\bottomrule
\endlastfoot
$\xi_\text{acute}$ (nm) & 0.425 & 0.065 & 0.303 & 0.541 & 293 \\
$\xi_\text{chronic}$ (nm) & 0.790 & 0.065 & 0.659 & 0.913 & 452 \\
$\xi_\text{healthy}$ (nm) & 0.797 & 0.048 & 0.717 & 0.887 & 325 \\
$\beta_\xi$ (protection exponent) & 2.33 & 0.51 & 1.49 & 3.26 & 251 \\
$\beta_\text{deloc}$ & 0.21 & 0.11 & 0.00 & 0.39 & 265 \\
NAA$_\text{base}$ & 1.12 & 0.06 & 1.02 & 1.23 & 418 \\
$k_\text{turnover}$ & 0.023 & 0.016 & 0.001 & 0.049 & 220 \\
$\sigma_\text{NAA}$ & 0.092 & 0.034 & 0.028 & 0.156 & 142 \\
$\sigma_\text{study}$ & 0.078 & 0.042 & 0.012 & 0.151 & 189 \\
era$_\text{effect}[0]$ & 0.012 & 0.031 & $-0.045$ & 0.068 & 312 \\
era$_\text{effect}[1]$ & $-0.008$ & 0.028 & $-0.058$ & 0.042 & 298 \\
\end{longtable}

\subsection{Enhanced Statistical Summary}

\begin{table}[htbp]
\centering
\caption{Enhanced statistical summary with Bayes Factors and effect sizes.}
\begin{tabular}{llr}
\toprule
\textbf{Category} & \textbf{Metric} & \textbf{Value} \\
\midrule
\multirow{3}{*}{Bayesian Evidence}
& $P(\xi_\text{acute} < \xi_\text{chronic})$ & $>0.999$ \\
& Bayes Factor (BF$_{10}$) & $>1000$ \\
& $\log(\text{BF}_{10})$ & 10.28 \\
\midrule
\multirow{4}{*}{Effect Sizes}
& Cohen's $d$ [95\% CI] & 5.63 [4.68, 6.58] \\
& Hedges' $g$ & 5.58 \\
& $\xi$ reduction (\%) & 46.2 \\
& $\xi$ reduction (nm) & 0.365 \\
\midrule
\multirow{4}{*}{Meta-Analysis}
& Pooled effect [95\% CI] & 0.099 [0.070, 0.128] \\
& $I^2$ (heterogeneity) & 0\% \\
& $\tau^2$ & 0.000 \\
& $Q$-statistic ($p$-value) & 0.60 ($p = 0.44$) \\
\midrule
\multirow{2}{*}{Model Selection}
& Full model weight & 0.89 \\
& $\Delta$WAIC (vs null) & 14.15 \\
\bottomrule
\end{tabular}
\end{table}

\clearpage
%%%%%%%%%%%%%%%%%%%%%%%%%%%%%%%%%%%%%%%%%%%%%%%%%%%%%%%%%%%%%%%%%%%%%%%%%%%%%%%
\section{Code Repository}
\label{sec:code}
%%%%%%%%%%%%%%%%%%%%%%%%%%%%%%%%%%%%%%%%%%%%%%%%%%%%%%%%%%%%%%%%%%%%%%%%%%%%%%%

All source code is available at: \url{https://github.com/Nyx-Dynamics/noise_decorrelation_hiv}. A frozen archive of the version used for this manuscript is deposited at Zenodo (DOI: \href{https://doi.org/10.5281/zenodo.18685010}{10.5281/zenodo.18685010}).

\subsection{Core Bayesian Model (v3.6)}

\textbf{File:} \texttt{quantum/bayesian\_v3\_6\_runner.py}

This module implements the primary hierarchical Bayesian model with ART-era effects.

\begin{lstlisting}[caption={Bayesian v3.6 Runner (Lines 1--100)}]
"""
Bayesian v3.6 runner with ART-era hierarchical effect while preserving
core mechanism Option A (xi -> Pi_xi).

- Loads bayesian_inputs_<ratio>.csv via data_loaders.load_bayesian_inputs
- Treats ART era as a categorical covariate (additive effect), not altering xi->Pi_xi
- Uses a simple forward mapping for NAA/Cr: baseline * Pi_xi_phase with study and era effects
- Retains validation-only fields in the annotated inputs written to outputs
- Writes per-run manifest with git/env details and checksums
"""
from __future__ import annotations

import argparse
from pathlib import Path
from typing import Dict

import arviz as az
import pandas as pd
import pymc as pm
import numpy as np

from quantum.data_loaders import load_bayesian_inputs
from quantum.utils.run_manifest import make_run_id, base_environment, write_manifest, Manifest

# Phase label normalization
_PHASE_NORM: Dict[str, str] = {
    'control': 'Control', 'healthy': 'Control', 'hc': 'Control', 'null': 'Control',
    'acute': 'Acute', 'acute_hiv': 'Acute',
    'chronic': 'Chronic', 'chronic_hiv': 'Chronic'
}


def norm_phase(val: str) -> str:
    if pd.isna(val):
        return 'Control'
    key = str(val).strip().lower()
    return _PHASE_NORM.get(key, val)


def parse_args():
    p = argparse.ArgumentParser(description='Bayesian v3.6 with ART-era hierarchical effect')
    p.add_argument('--ratio', required=True, choices=['3_1_1', '1_1_1', '1_2_1', '2_1'])
    p.add_argument('--era', default='both', choices=['pre_modern', 'post_modern', 'both'])
    p.add_argument('--draws', type=int, default=1500)
    p.add_argument('--tune', type=int, default=1000)
    p.add_argument('--chains', type=int, default=4)
    p.add_argument('--target-accept', type=float, default=0.99)
    p.add_argument('--seed', type=int, default=2025)
    p.add_argument('--output-root', default='results/bayesian_v3_6')
    p.add_argument('--run-notes', default='')
    # Mechanism parameters
    p.add_argument('--beta-xi-mean', type=float, default=1.89)
    p.add_argument('--beta-xi-sd', type=float, default=0.25)
    p.add_argument('--xi-baseline-nm', type=float, default=0.66)
    p.add_argument('--baseline-naa-cr', type=float, default=2.2)
    return p.parse_args()


def prepare(df: pd.DataFrame, era: str) -> pd.DataFrame:
    df = df.copy()
    # Filter era if requested
    if era != 'both' and 'art_era' in df.columns:
        df = df[df['art_era'] == era]

    # Normalize phases
    if 'Phase' in df.columns:
        df['Phase'] = df['Phase'].apply(norm_phase)
    else:
        raise ValueError("bayesian_inputs CSV must include 'Phase' column")

    # Auto-detect schema: filter to NAA/Cr metabolite if present
    if 'Metabolite' in df.columns:
        meta_norm = df['Metabolite'].astype(str).str.replace(' ', '', regex=False).str.upper()
        naa_mask = meta_norm.eq('NAA/CR') | meta_norm.eq('NAA/CRR') | meta_norm.eq('NAA/CRATIO')
        if naa_mask.any():
            df = df[naa_mask].copy()

    # Rename columns if needed
    if ('NAA_mean' not in df.columns or 'NAA_SE' not in df.columns) and \
       {'Mean', 'SE'}.issubset(df.columns):
        df = df.rename(columns={'Mean': 'NAA_mean', 'SE': 'NAA_SE'})

    # Require NAA mean/SE
    needed = ['NAA_mean', 'NAA_SE']
    missing_cols = [c for c in needed if c not in df.columns]
    if missing_cols:
        raise ValueError(f"Missing required columns: {missing_cols}")
    df = df[~df['NAA_mean'].isna() & ~df['NAA_SE'].isna()].copy()

    # Study id for random effects
    if 'study' not in df.columns:
        if 'Study' in df.columns:
            df['study'] = df['Study']
        else:
            df['study'] = 'pooled'
    df['study_id'] = df['study'].astype('category').cat.codes

    return df
\end{lstlisting}

\begin{lstlisting}[caption={Bayesian v3.6 Runner - Model Building (Lines 127--175)}]
def build_model(df: pd.DataFrame, beta_xi_mean: float, beta_xi_sd: float,
                xi_baseline_nm: float, baseline_naacr: float):
    # Indexing helpers
    phases = ['Control', 'Acute', 'Chronic']
    phase_idx = df['Phase'].astype('category').cat.set_categories(phases).cat.codes.values
    study_ids = df['study_id'].values
    era_idx = df['art_era_idx'].values

    N = len(df)
    n_studies = int(df['study_id'].nunique())

    with pm.Model() as model:
        # Mechanism: beta_xi prior
        beta_xi = pm.TruncatedNormal('beta_xi', mu=beta_xi_mean, sigma=beta_xi_sd, lower=0.0)

        # Latent xi per phase (log-parameterization for improved geometry)
        log_xi_ctrl = pm.Normal('log_xi_control', mu=np.log(0.66), sigma=0.15)
        log_xi_acute = pm.Normal('log_xi_acute', mu=np.log(0.61), sigma=0.15)
        log_xi_chron = pm.Normal('log_xi_chronic', mu=np.log(0.80), sigma=0.15)
        xi_nm_ctrl = pm.Deterministic('xi_nm_control', pm.math.exp(log_xi_ctrl))
        xi_nm_acute = pm.Deterministic('xi_nm_acute', pm.math.exp(log_xi_acute))
        xi_nm_chron = pm.Deterministic('xi_nm_chronic', pm.math.exp(log_xi_chron))
        log_xi = pm.math.stack([log_xi_ctrl, log_xi_acute, log_xi_chron])

        # Protection factor per phase: (xi_ref/xi)^{beta_xi}
        Pi_xi_phase = pm.math.exp((-beta_xi) * (log_xi - np.log(xi_baseline_nm)))

        # Era additive effect
        era_effect = pm.Normal('era_effect', mu=0.0, sigma=0.10, shape=2)

        # Study random effects
        study_sd = pm.HalfNormal('study_sd', sigma=0.10)
        study_offset = pm.Normal('study_offset', mu=0.0, sigma=1.0, shape=n_studies)
        study_effect = pm.Deterministic('study_effect', study_offset * study_sd)

        # Observation noise
        se_scale = pm.HalfNormal('se_scale', sigma=0.20)

        # Expected mean per observation
        mu_base = baseline_naacr * Pi_xi_phase[phase_idx]
        era_term = pm.math.switch(pm.math.eq(era_idx, -1), 0.0, era_effect[era_idx])
        mu = mu_base * (1.0 + era_term + study_effect[study_ids])

        # Likelihood
        sigma = pm.math.maximum(1e-6, se_scale * df['NAA_SE'].values)
        pm.Normal('naa_obs', mu=mu, sigma=sigma, observed=df['NAA_mean'].values)

    return model
\end{lstlisting}

\subsection{Enzyme Kinetics Validation Model (v4.0)}

\textbf{File:} \texttt{quantum/bayesian\_enzyme\_v4.py}

This module implements the mechanistic enzyme kinetics model for validation.

\begin{lstlisting}[caption={Enzyme Kinetics Model - Forward Model (Lines 77--166)}]
def forward_model_enzyme(xi_acute, xi_chronic, beta_xi, gamma_coh,
                        viral_damage_acute, viral_damage_chronic,
                        membrane_acute, membrane_chronic,
                        coh_acute=0.95, coh_chronic=0.80):
    """
    Forward model using enzyme kinetics.

    Parameters
    ----------
    xi_acute : float
        Correlation length in acute HIV (m)
    xi_chronic : float
        Correlation length in chronic HIV (m)
    beta_xi : float
        Protection factor exponent
    gamma_coh : float
        Coherence coupling exponent
    viral_damage_acute : float
        Viral damage factor in acute (0-1)
    viral_damage_chronic : float
        Viral damage factor in chronic (0-1)
    membrane_acute : float
        Membrane turnover in acute (>1 = elevated)
    membrane_chronic : float
        Membrane turnover in chronic (>1 = elevated)

    Returns
    -------
    NAA_pred : array
        [NAA_healthy, NAA_acute, NAA_chronic] in MRS units
    Cho_pred : array
        [Cho_healthy, Cho_acute, Cho_chronic] in MRS units
    """

    # Healthy baseline
    xi_healthy = 0.75e-9  # Reference
    Pi_healthy = compute_protection_factor(xi_healthy, beta_xi=beta_xi)
    eta_healthy = coherence_modulation(0.85, gamma=gamma_coh)

    enzymes_healthy = EnzymeKinetics(
        Pi_xi=Pi_healthy,
        eta_coh=eta_healthy,
        viral_damage_factor=1.0
    )
    NAA_h, Cho_h = enzymes_healthy.integrate(
        duration_days=60,
        membrane_turnover=1.0
    )

    # Acute HIV
    Pi_acute = compute_protection_factor(xi_acute, beta_xi=beta_xi)
    eta_acute = coherence_modulation(coh_acute, gamma=gamma_coh)

    enzymes_acute = EnzymeKinetics(
        Pi_xi=Pi_acute,
        eta_coh=eta_acute,
        viral_damage_factor=viral_damage_acute
    )
    NAA_a, Cho_a = enzymes_acute.integrate(
        duration_days=60,
        membrane_turnover=membrane_acute
    )

    # Chronic HIV
    Pi_chronic = compute_protection_factor(xi_chronic, beta_xi=beta_xi)
    eta_chronic = coherence_modulation(coh_chronic, gamma=gamma_coh)

    enzymes_chronic = EnzymeKinetics(
        Pi_xi=Pi_chronic,
        eta_coh=eta_chronic,
        viral_damage_factor=viral_damage_chronic
    )
    NAA_c, Cho_c = enzymes_chronic.integrate(
        duration_days=60,
        membrane_turnover=membrane_chronic
    )

    # Convert from molar to MRS units (relative to creatine)
    creatine = 8.0e-3

    NAA_pred = np.array([NAA_h, NAA_a, NAA_c]) / creatine
    Cho_pred = np.array([Cho_h, Cho_a, Cho_c]) / creatine

    return NAA_pred, Cho_pred
\end{lstlisting}

\begin{lstlisting}[caption={Enzyme Kinetics Model - PyMC Model (Lines 172--314)}]
def build_enzyme_model():
    """
    Build PyMC model with enzyme kinetics.

    KEY PARAMETERS:
    - xi_acute, xi_chronic: Noise correlation lengths
    - beta_xi: Protection factor exponent (expect ~2)
    - gamma_coh: Coherence coupling exponent
    - viral_damage_*: Direct damage to enzymes
    - membrane_*: Membrane turnover rates
    """

    with pm.Model() as model:

        # PRIORS
        # Noise correlation lengths (informed by v3.6)
        xi_acute = pm.TruncatedNormal(
            'xi_acute',
            mu=0.50e-9,
            sigma=0.15e-9,
            lower=0.35e-9,
            upper=0.70e-9
        )

        xi_chronic = pm.TruncatedNormal(
            'xi_chronic',
            mu=0.78e-9,
            sigma=0.10e-9,
            lower=0.70e-9,
            upper=0.90e-9
        )

        # Protection factor exponent (informed by v3.6: beta = 1.731)
        beta_xi = pm.TruncatedNormal(
            'beta_xi',
            mu=1.75,
            sigma=0.50,
            lower=0.5,
            upper=3.5
        )

        # Coherence coupling exponent
        gamma_coh = pm.TruncatedNormal(
            'gamma_coh',
            mu=1.5,
            sigma=0.50,
            lower=0.5,
            upper=3.0
        )

        # Viral damage factors (0-1, 1 = no damage)
        viral_damage_acute = pm.Beta('viral_damage_acute', alpha=19, beta=1)
        viral_damage_chronic = pm.Beta('viral_damage_chronic', alpha=9, beta=1)

        # Membrane turnover
        membrane_acute = pm.TruncatedNormal('membrane_acute', mu=2.0, sigma=0.5,
                                           lower=1.0, upper=4.0)
        membrane_chronic = pm.TruncatedNormal('membrane_chronic', mu=1.2, sigma=0.3,
                                             lower=1.0, upper=2.0)

        # FORWARD MODEL
        NAA_pred, Cho_pred = forward_model_enzyme_op(
            xi_acute, xi_chronic, beta_xi, gamma_coh,
            viral_damage_acute, viral_damage_chronic,
            membrane_acute, membrane_chronic
        )

        # LIKELIHOOD
        sigma_NAA = pm.HalfNormal('sigma_NAA', sigma=0.06)
        sigma_Cho = pm.HalfNormal('sigma_Cho', sigma=0.03)

        NAA_likelihood = pm.Normal('NAA_obs', mu=NAA_pred, sigma=sigma_NAA,
                                   observed=NAA_OBS)
        Cho_likelihood = pm.Normal('Cho_obs', mu=Cho_pred, sigma=sigma_Cho,
                                   observed=CHO_OBS)

        # DERIVED QUANTITIES
        Pi_acute = pm.Deterministic('Pi_acute', (0.8e-9 / xi_acute) ** beta_xi)
        Pi_chronic = pm.Deterministic('Pi_chronic', (0.8e-9 / xi_chronic) ** beta_xi)
        protection_ratio = pm.Deterministic('protection_ratio', Pi_acute / Pi_chronic)

    return model
\end{lstlisting}

\subsection{Statistical Improvements Module}

\textbf{File:} \texttt{scripts/statistical\_improvements.py}

This module implements enhanced statistical reporting including Bayes Factors, effect sizes, and meta-analysis.

\begin{lstlisting}[caption={Statistical Improvements - Bayes Factor (Lines 145--196)}]
def calculate_bayes_factor_savage_dickey(
    posterior_mean: float,
    posterior_sd: float,
    prior_mean: float,
    prior_sd: float,
    null_value: float = 0.0
) -> BayesFactorResult:
    """
    Calculate Bayes Factor using Savage-Dickey density ratio.

    BF_10 = P(theta=theta_0 | prior) / P(theta=theta_0 | posterior)

    For testing H0: xi_acute >= xi_chronic (no protection)
    vs H1: xi_acute < xi_chronic (protective effect)
    """
    # Prior density at null
    prior_density = stats.norm.pdf(null_value, prior_mean, prior_sd)

    # Posterior density at null
    posterior_density = stats.norm.pdf(null_value, posterior_mean, posterior_sd)

    # Avoid division by zero
    if posterior_density < 1e-300:
        bf_10 = np.inf
        log_bf_10 = np.inf
    else:
        bf_10 = prior_density / posterior_density
        log_bf_10 = np.log(prior_density) - np.log(posterior_density)

    # Interpretation (Kass & Raftery 1995)
    if bf_10 > 100:
        interpretation = "Decisive evidence for H1"
    elif bf_10 > 30:
        interpretation = "Very strong evidence for H1"
    elif bf_10 > 10:
        interpretation = "Strong evidence for H1"
    elif bf_10 > 3:
        interpretation = "Substantial evidence for H1"
    elif bf_10 > 1:
        interpretation = "Weak evidence for H1"
    else:
        interpretation = "Evidence favors H0"

    return BayesFactorResult(
        bf_10=bf_10,
        log_bf_10=log_bf_10,
        interpretation=interpretation,
        h0_description="xi_acute >= xi_chronic (no protective effect)",
        h1_description="xi_acute < xi_chronic (protective paradox)",
        method="Savage-Dickey density ratio"
    )
\end{lstlisting}

\begin{lstlisting}[caption={Statistical Improvements - Meta-Analysis (Lines 334--411)}]
def random_effects_meta_analysis(
    effects: np.ndarray,
    variances: np.ndarray,
    method: str = 'DL'
) -> MetaAnalysisResult:
    """
    Perform random-effects meta-analysis using DerSimonian-Laird method.

    Parameters
    ----------
    effects : array
        Effect sizes from each study
    variances : array
        Variance of effect sizes
    method : str
        'DL' for DerSimonian-Laird, 'REML' for restricted maximum likelihood

    Returns
    -------
    MetaAnalysisResult with pooled effect, heterogeneity statistics
    """
    k = len(effects)

    if k < 2:
        return MetaAnalysisResult(
            pooled_effect=effects[0] if k == 1 else np.nan,
            pooled_se=np.sqrt(variances[0]) if k == 1 else np.nan,
            pooled_ci=(np.nan, np.nan),
            q_statistic=np.nan,
            q_pvalue=np.nan,
            i_squared=0.0,
            tau_squared=0.0,
            n_studies=k,
            method=method
        )

    # Fixed-effects weights
    w = 1 / variances

    # Fixed-effects pooled estimate
    theta_fe = np.sum(w * effects) / np.sum(w)

    # Q statistic (heterogeneity test)
    Q = np.sum(w * (effects - theta_fe)**2)
    df = k - 1
    q_pvalue = 1 - stats.chi2.cdf(Q, df)

    # DerSimonian-Laird estimate of tau^2
    c = np.sum(w) - np.sum(w**2) / np.sum(w)
    tau2 = max(0, (Q - df) / c)

    # Random-effects weights
    w_re = 1 / (variances + tau2)

    # Random-effects pooled estimate
    theta_re = np.sum(w_re * effects) / np.sum(w_re)
    se_re = np.sqrt(1 / np.sum(w_re))

    # 95% CI
    ci = (theta_re - 1.96*se_re, theta_re + 1.96*se_re)

    # I^2 statistic
    if Q > 0:
        i_squared = max(0, ((Q - df) / Q) * 100)
    else:
        i_squared = 0.0

    return MetaAnalysisResult(
        pooled_effect=theta_re,
        pooled_se=se_re,
        pooled_ci=ci,
        q_statistic=Q,
        q_pvalue=q_pvalue,
        i_squared=i_squared,
        tau_squared=tau2,
        n_studies=k,
        method=method
    )
\end{lstlisting}

\subsection{Data Loader Utilities}

\textbf{File:} \texttt{quantum/data\_loaders.py}

This module handles data loading with ART-era classification and mechanistic parameter derivation.

\begin{lstlisting}[caption={Data Loaders (Complete File)}]
"""
Data loaders for ratio-comparison datasets with ART-era classification and
mechanistic Option A (xi -> Pi_xi) derivation.
"""
from __future__ import annotations

from pathlib import Path
from typing import Literal, Tuple

import numpy as np
import pandas as pd

RATIO_DIR = Path('data/extracted_expanded/data_ratios_comparison')

Era = Literal['pre_modern', 'post_modern', 'unknown']

PAPER_YEAR: dict[str, int] = {
    'Valcour_2015': 2015,
    'Young_2014': 2014,
    'Sailasuta_2012': 2012,
    'Chang_2002': 2002,
}

ART_CUTOFF_YEAR = 2006  # <= pre_modern, >=2007 post_modern


def classify_era(year: float | int | None, cutoff: int = ART_CUTOFF_YEAR) -> Era:
    if pd.isna(year):
        return 'unknown'
    try:
        y = int(year)
    except Exception:
        return 'unknown'
    return 'pre_modern' if y <= cutoff else 'post_modern'


def attach_era(df: pd.DataFrame, cutoff: int = ART_CUTOFF_YEAR) -> pd.DataFrame:
    df = df.copy()
    year_series = None
    for col in ['measurement_year', 'scan_year', 'publication_year']:
        if col in df.columns:
            year_series = df[col]
            break

    if year_series is None:
        if 'study' in df.columns:
            df['publication_year'] = df['study'].map(PAPER_YEAR)
            year_series = df['publication_year']
        else:
            df['publication_year'] = np.nan
            year_series = df['publication_year']

    df['art_era'] = [classify_era(y, cutoff) for y in year_series]
    df['art_era_idx'] = df['art_era'].map({'pre_modern': 0, 'post_modern': 1, 'unknown': -1})
    return df


def add_derived_covariates(df: pd.DataFrame) -> pd.DataFrame:
    """Compute derived covariates when inputs exist."""
    df = df.copy()
    if 'pVL' in df.columns:
        with np.errstate(divide='ignore'):
            df['log_VL'] = np.log10(pd.to_numeric(df['pVL'], errors='coerce').clip(lower=1.0))
    if {'CD4', 'CD8'}.issubset(df.columns):
        with np.errstate(divide='ignore', invalid='ignore'):
            c4 = pd.to_numeric(df['CD4'], errors='coerce')
            c8 = pd.to_numeric(df['CD8'], errors='coerce')
            df['CD4_CD8_ratio'] = c4 / c8
    return df


def compute_Pi_xi_from_xi_nm(xi_nm: pd.Series, beta_xi: float = 1.89,
                            xi_baseline_nm: float = 0.66) -> pd.Series:
    xi = pd.to_numeric(xi_nm, errors='coerce').astype(float).clip(lower=1e-9)
    return (xi / float(xi_baseline_nm)) ** (-beta_xi)


def load_enzyme_inputs(ratio: str, beta_xi: float = 1.89,
                      xi_baseline_nm: float = 0.66) -> Tuple[pd.DataFrame, Path]:
    path = RATIO_DIR / f'enzyme_inputs_{ratio}.csv'
    if not path.exists():
        raise FileNotFoundError(f"Missing enzyme inputs for ratio '{ratio}': {path}")
    df = pd.read_csv(path)
    df = attach_era(df)
    if 'xi_estimate_nm' in df.columns:
        df['Pi_xi'] = compute_Pi_xi_from_xi_nm(df['xi_estimate_nm'],
                                               beta_xi=beta_xi,
                                               xi_baseline_nm=xi_baseline_nm)
        df['computed_Pi_xi'] = df['Pi_xi']
    else:
        raise ValueError("enzyme_inputs CSV must include 'xi_estimate_nm' column")
    df = add_derived_covariates(df)
    return df, path


def load_bayesian_inputs(ratio: str) -> Tuple[pd.DataFrame, Path]:
    path = RATIO_DIR / f'bayesian_inputs_{ratio}.csv'
    if not path.exists():
        raise FileNotFoundError(f"Missing bayesian inputs for ratio '{ratio}': {path}")
    df = pd.read_csv(path)
    df = attach_era(df)
    df = add_derived_covariates(df)
    return df, path
\end{lstlisting}

\newpage
\subsection{Software Environment}

\begin{table}[htbp]
\centering
\caption{Software dependencies and versions used in this analysis.}
\begin{tabular}{llr}
\toprule
\textbf{Package} & \textbf{Purpose} & \textbf{Version} \\
\midrule
Python & Programming language & 3.11+ \\
PyMC & Bayesian inference & 5.x \\
ArviZ & Posterior analysis & 0.18+ \\
NumPy & Numerical computing & 1.24+ \\
Pandas & Data manipulation & 2.0+ \\
SciPy & Statistical functions & 1.11+ \\
Matplotlib & Visualization & 3.8+ \\
scikit-learn & Machine learning utilities & 1.5.0+ \\
tqdm & Progress bars & 4.66.3+ \\
\bottomrule
\end{tabular}
\end{table}

\subsection{Reproducibility}

All analysis code, data, and results are archived at Zenodo (DOI: \href{https://doi.org/10.5281/zenodo.18685010}{10.5281/zenodo.18685010}) and maintained at \url{https://github.com/Nyx-Dynamics/noise\_decorrelation\_hiv}.

To reproduce the complete analysis pipeline:

\begin{enumerate}
    \item Clone the repository: \\
    \texttt{git clone https://github.com/Nyx-Dynamics/noise\_decorrelation\_hiv.git} \\
    \texttt{cd noise\_decorrelation\_hiv}

    \item Create a virtual environment and install dependencies: \\
    \texttt{python3 -m venv .venv \&\& source .venv/bin/activate} \\
    \texttt{pip install -r requirements.txt}

    \item Reproduce all results (models v3.6, v1.0, v2.0, v4.0): \\
    \texttt{python reproduce\_all.py} \\
    or equivalently: \texttt{make reproduce}
\end{enumerate}

All random seeds are fixed for deterministic reproduction. The pipeline runs four convergent Bayesian models, generates all manuscript figures, and produces cryptographic checksums for output verification. Expected runtime on a modern workstation (8+ cores, 16 GB RAM): approximately 30--60 minutes for full MCMC sampling across all models. Individual models can also be run independently (see \texttt{QUICKSTART.md}).

\newpage
%%%%%%%%%%%%%%%%%%%%%%%%%%%%%%%%%%%%%%%%%%%%%%%%%%%%%%%%%%%%%%%%%%%%%%%%%%%%%%%
\section*{References}
%%%%%%%%%%%%%%%%%%%%%%%%%%%%%%%%%%%%%%%%%%%%%%%%%%%%%%%%%%%%%%%%%%%%%%%%%%%%%%%

\begin{enumerate}
    \item Kass RE, Raftery AE. Bayes Factors. \textit{J Am Stat Assoc}. 1995;90(430):773-795.
    \item DerSimonian R, Laird N. Meta-analysis in clinical trials. \textit{Control Clin Trials}. 1986;7(3):177-188.
    \item Watanabe S. Asymptotic equivalence of Bayes cross validation and widely applicable information criterion in singular learning theory. \textit{J Mach Learn Res}. 2010;11:3571-3594.
    \item Gelman A, Carlin JB, Stern HS, Dunson DB, Vehtari A, Rubin DB. \textit{Bayesian Data Analysis}. 3rd ed. CRC Press; 2013.
    \item Vehtari A, Gelman A, Gabry J. Practical Bayesian model evaluation using leave-one-out cross-validation and WAIC. \textit{Stat Comput}. 2017;27(5):1413-1432.
\end{enumerate}

\end{document}
